\documentclass{article}

% Install the official ICLR 2027 style beside this file before submission.
\usepackage{iclr2027_conference,times}
\usepackage{amsmath,amssymb,booktabs,graphicx,microtype,url}

\title{Experience Models: Per-Life Parametric Learning for Acting Agents}
\author{Anonymous authors}

\newcommand{\tbd}{\textbf{[TBD]}}
\newcommand{\unproven}[1]{\textbf{[UNPROVEN: #1]}}

\begin{document}
\maketitle

\begin{abstract}
Agent post-training teaches a model how to act before deployment, but does not
teach one deployed agent to improve from its own subsequent life. We study
whether process-level parenting can bootstrap that ability. One parent teaches
one child how to think through target-blind practice tasks and
thought-to-action correction, then disappears. The child enters a fresh
compiler-optimization gym with a frozen base model, ordinary external agent
resources, and a per-life low-rank state updated by a Think--Dream--Sleep
loop. A matched $2\times2$ design crosses parenting with deployment-time
parametric consolidation, separating inherited starting competence, online
learning, and their interaction; a fit-free frozen-parameter agent with the
same validated evolving textual memory supplies the paper-facing
frozen-parameter active-memory reference. We measure held-out value as a function of
lifetime under equal generated-token budgets, with proposal/routing assays,
forward-transfer and retention panels, adapter-removal controls, and complete
resource accounting. \unproven{Replace this sentence with the registered
parent-absent learning result, or the strongest surviving lower-rung result,
only after the sealed run.}
\end{abstract}

\section{Introduction}

Pretraining and fleet-level agent post-training amortize experience across a
population before deployment. External memory lets an individual retain a
life while leaving its acting policy fixed. Our question is complementary:
can one deployed agent convert its own action--outcome experience into a
persistent parametric change that improves how it acts later?

The difficulty is not merely storing facts. The writer must preserve the
agent's executable interface, distinguish local episodes from reusable
structure, avoid reinforcing its own formatting accidents, and remain stable
over repeated writes. These requirements turn personal experiential learning
into an online stability--plasticity problem.

Our intended contributions are:
\begin{enumerate}
  \item A one-parent/one-child protocol that teaches process-level thinking
        through action and public consequence, while keeping the parent out
        of the verifier role and removing all parental text at deployment.
  \item A prospective $2\times2$ test of parenting by deployment-time
        parametric consolidation, with a common development-gated active-text
        memory system and separate measures of starting competence and
        learning rate.
  \item A provenance-bound Think--Dream--Sleep architecture with separate
        measures of proposal content, executable routing, forward transfer,
        backward retention, and writer-induced failure.
  \item \unproven{A registered positive parent-absent lifetime result and,
        conditionally, a terminal parametric-carrier mechanism result.}
\end{enumerate}

\section{Setting and hypotheses}

Let a life produce an append-only stream
\(e_t=(c_t,z_t,a_t,o_t)\), where \(c_t\) is rendered state, \(z_t\) is
free-form thinking, \(a_t\) is a typed action, and \(o_t\) is a public
outcome. A frozen base policy \(\pi_0\) acts with a per-life low-rank state
\(\Delta_t\). Periodically, an experience compiler admits supported native
continuations and produces \(\Delta_{t+1}\).

The primary hypothesis is that process-only parenting increases the benefit
of deployment-time per-life consolidation after the parent is removed. The
primary estimand is therefore the interaction between parenting and personal
writes over the held-out lifetime curve, not the final score of a single
parented agent. Secondary hypotheses test parent-absent starting competence,
unparented online learning, forward transfer, backward retention, and whether
any gain survives adapter-removal and carrier interventions.

``Model meta-intelligence'' is used only if prior target-blind parenting
improves the rate at which the parent-absent child learns a fresh environment.
Continual weight change alone, or a higher score at deployment entry, does not
establish this claim.

\section{Architecture}

\paragraph{Think.}
One model generates a continuous reasoning stream and typed tool calls.
Thinking, planning, reflection, and context reconciliation are uses of the
same inference operation under different admitted state; they are not
separately trained modules.

\paragraph{Ledger and conscious state.}
Every invocation, native response envelope, dispatched action, public
outcome, correction, and compiled row is recorded prospectively. Context may
be distilled only after a validated successor state is published; the ledger
keeps the operation reversible.

\paragraph{Sleep-write.}
Sleep is the write boundary, not a second intelligence. The compiler emits
native, response-only training targets whose action bytes actually dispatched
and whose support is bound to public outcomes. Context, tool results, parent
text, and old history are masked inputs. The adapter is trained from the
frozen base plus the cumulative accepted corpus.

\paragraph{Strong active-text control.}
The proposed frozen-parameter opponent is
\texttt{ACTIVE\_TEXT\_NATIVE-v2.1-AUTO}: an isolated append-only public ledger
plus a deterministic lexical and witnessed-transition-graph retriever. It has
no reflector, curator, model-generated query, learned updater, hidden graph,
or parameter change. Before each of the same first 16 actor calls, retrieval
is driven automatically by the public objective, current state, and the
child's prior task-local thoughts; returned records count against a fixed
memory-token budget. This makes the reference an adapting memory \emph{system}
with a frozen actor, and asks whether personal LoRA writes add value beyond a
strong textual path. The baseline remains a proposal until its exact source,
use, leakage, headroom, and doubled-bandwidth certificates pass.

\paragraph{Parenting.}
A parent may propose process-level credit assignment but never supplies a
hidden answer. A correction is admitted only after the child applies it on a
separate training situation and the public environment supports it. Parent
scaffolds are faded in the write, and all headline probes remove parent text.
The paper uses one parent and one child; classroom, multi-teacher, peer, and
population mechanisms are outside scope.

\section{Benchmark and protocol}

\subsection{Controlled experiential environment}
\tbd{} Specify the target-blind nursery tasks used for process correction and
the conditional terminal carrier assay. Nursery content must exclude
CompilerGym programs, LLVM pass sequences, benchmark outcomes, and probe
targets.

\subsection{Compiler optimization transfer environment}
We use a deterministic CompilerGym objective as an externally grounded
transfer demonstration. Programs, exact action tuples, and content hashes are
partitioned before target enumeration. Evaluation uses equal generated-token
budgets, common random generation, isolated ledgers, and fresh environment
resets.

\subsection{Parenting and deployment transfer}

One parent teaches one child process-level thinking through target-blind
practice tasks. The child must turn corrections into actions on separate
training situations; the public environment admits supported corrections into
sleep. At deployment, the parent, correction text, nursery transcript, and
practice ledger are absent. The child retains only the permitted learned
parametric state and enters a fresh task stream.

The headline compares this parented Think--Dream--Sleep agent with a fit-free
frozen-parameter active-memory agent under identical deployment affordances and generated-token budgets. A
$2\times2$ diagnostic crosses parenting with deployment-time consolidation:
the parented checkpoint and a dose-matched unparented practice-written
checkpoint, each forked with deployment writes off/on. All services own an
isolated, frozen-policy ACE/ReMem-style playbook updater and pinned hybrid
retriever; the common mechanism must pass prospective write, read/use,
end-to-end, and task-headroom gates before it is called a strong baseline.
This separates inherited starting competence from a higher learning rate and
from the interaction of parenting with the per-life learning loop. No
classroom or population mechanism is used.

Let $Y_{p,\ell}(t)$ denote held-out value after deployment lifetime $t$, where
$p$ indicates parenting and $\ell$ permits deployment-time consolidation. The
parenting contribution to online learning is the predeclared
difference-in-differences
\[
  I(t)=\bigl[Y_{1,1}(t)-Y_{1,0}(t)\bigr]
       -\bigl[Y_{0,1}(t)-Y_{0,0}(t)\bigr].
\]
We also compare change from each cell's pre-deployment checkpoint and fit the
registered life-level learning curve. A higher parented starting score without
positive $I(t)$ is inherited competence, not evidence that parenting taught
the child to learn better.

If only one child is parented, we report this extension as a mechanistic case
study; deployment tasks do not replicate its childhood. A general parenting
effect requires independently raised one-on-one child lives, all taught by the
same parent under the frozen protocol.

\subsection{Causal arms and matched resources}

The four deployment cells are unparented frozen (\texttt{U0}), parented frozen
(\texttt{P0}), unparented continual (\texttt{U1}), and parented continual
(\texttt{P1}). Both unparented cells inherit a target-blind, dose-matched
childhood LoRA written only from their own supported practice, so \texttt{U0}
is not the pristine base. The fit-free reference \texttt{R0} retains raw
pinned actor weights while its isolated active-text playbook evolves, so the
complete system is frozen-parameter rather than memory-static.
The paper-facing two-agent comparison is \texttt{P1} versus \texttt{R0}; the
P/U $2\times2$ identifies the parenting interaction.

All cells receive matched tools, base model, generated-token budgets, task
opportunities, public outcomes, files, and skills. Longitudinal systems are
constructed from the same branch-local public event stream; the parametric
arms compile it into LoRA state, while the fit-free reference exposes it only
through \texttt{ACTIVE\_TEXT\_NATIVE-v2.1-AUTO}. Only parenting and/or per-life
parametric writes are disabled in the registered $2\times2$. The hierarchy tests the parenting-
by-write interaction, then \texttt{P1}-\texttt{P0} as the total closed-loop
benefit of LoRA writes on top of strong active text, then
\texttt{P1}-\texttt{R0} as the full-system public contrast. Raw-trajectory
distillation, a bounded descriptive-only terminal LEAFE-style system, binding
derangement, and adapter removal are staged diagnostics rather than coequal
longitudinal arms. The LEAFE-style result supports no powered superiority
claim.
Accepted rows, supervised target tokens, optimization steps, model calls,
wall time, and GPU time are reported.

\subsection{Endpoints and statistics}

Primary endpoints are held-out task value and normalized area under the
value--lifetime curve, analyzed through the registered $2\times2$ interaction
and change from deployment entry. We additionally report early learning
slope, value per generated token, free typed-call compliance, typed-forced
proposal quality, forward transfer, backward retention, and forgetting. The
independent unit is a separately raised child life; checkpoints, programs,
actions, and decoding samples are repeated measures. \tbd{} Insert the frozen
parent identity/policy, parenting-cycle count, preregistered SESOIs, blinded
variance rule, intervals, root count, missing-action value, and joint success
rule.

\section{Results}

\subsection{Exploratory failure: useful proposals leave the action channel}

In one exploratory life, the terminal periodic adapter scored $0.0511$ under
the registered strict parser, versus $0.4859$ with that checkpoint's adapter
removed. A post-hoc permissive parser recovered a first proposed action on all
eight saved tasks; executing those exact saved strings once scored $0.5287$.
This localizes a writer-induced routing failure but is not a learning benefit:
generation was unseeded, context policy was confounded, and the permissive
analysis was defined after observing the failure.

\subsection{DEV-only worksheet prompt-package robustness (SEQ192)}

On September 13, 2026, three existing SEQ142 prediction adapters were tested
without new fits or updates against disabled-adapter OFF on 24 shared fresh
authored cases, each in FULL and MINIMAL views. Each view contains six cases
per supported-true, supported-false, missing, and conflicting evidence group.
Primary typed-content counts retain malformed and truncated outputs.

\begin{center}
\begin{tabular}{lrrrrr}
\toprule
Seed & OFF FULL & Post FULL & OFF MINIMAL & Post MINIMAL & $\Delta$ MINIMAL \\
\midrule
0 & 16 & 24 & 14 & 24 & $+10$ \\
1 & 16 & 23 & 14 & 1 & $-13$ \\
2 & 16 & 19 & 14 & 18 & $+4$ \\
\bottomrule
\end{tabular}
\end{center}
Every count has denominator 24. The predeclared continuation rule is
\textbf{FALSE}: positive MINIMAL contrast and post content at least 20/24
were required in every learner; seeds 1 and 2 fail. This branch stops without
prompt, dose, scorer, or token-cap rescue. FULL improves in all learners, but
uniform prompt-package robustness fails. Strict MINIMAL post counts are
0/0/18 versus OFF 0/0/0: seed 0's 24 content-correct answers are all fenced;
seed 1 has 23 unparseable outputs and one fenced answer, with one length
termination and 23 stop terminations. Content and strict format are distinct.

The original SEQ142 selected action-ID residues encode target labels; its
48/48 prediction result therefore does not establish evidence reading.
SEQ192 uses fresh disjoint triples and 24 distinct selected residues but
\textbf{has no same-ID counterfactual test}. FULL/MINIMAL differ in packaging,
so this tests prompt-package robustness only, not isolated conditional
evidence use or exclusion of all ID shortcuts. It supplies no matched-trained
parenting comparison, retention test, independent fact-bank replication,
new-family transfer, general H1/H2 or P1 result, or mechanism-freeze claim.

Primary attempts 2/2/1 for seeds 0/1/2 were selected before outcome inspection.
Original seed 0/1 attempts remain ineligible after controller custody-check
failures, not zero-score observations or repaired receipts. Resource accounting
is $288$ eligible $+192$ failed-attempt $=480$ actual calls, with zero fits,
updates, or parent calls. Shared cases and repeated frozen OFF do not supply
independent banks. Retention-v2 runtime integration is separate and supplies
no result here; the manuscript and research mission remain incomplete.

The result memo and durable audit JSON are repository-relative sources:
\begin{quote}\small
\path{research_notes/analysis/2026-09-13_level1_prediction_transfer_result.md}\\
\path{research_notes/astra_memos/receipts_20260912/astra_prediction_transfer_audit_20260913_result.json}\\
Audit JSON SHA256:
\path{9ca1433d5cce275cbc9ae8af1a36a43e7f93351d28eb823e9d6e5cde8c821344}.
\end{quote}

\subsection{Bounded replay retention on child EVENT records (SEQ195)}

We tested sequential writing on one exposed DEV bank containing four A and
four B EVENT records. The three optimizer seeds used the same semantic bank
and therefore test fit repeatability, not independent-bank reliability. Every
condition was read cold in both W0 and W8 views. In all three seeds and both
views, the initial A write moved A from $0/4$ to $4/4$. A subsequent B-only
write learned B ($4/4$) but forgot A ($0/4$), whether matched to the new dose
or to replay's total work. Replaying A while writing B retained A and learned
B ($4/4$ and $4/4$); clean cumulative retraining reached the same descriptive
endpoint.

The registered replay-minus-new-only A contrast was $+12/12$ pooled learner
cells at each view, with no B loss. All 288 accepted readout calls were
exact-hash audited and had zero truncations. Across eligible and failed
attempts, the campaign used 20 fits, 6,400 updates, and 25,600 presentations.
No single control simultaneously matches replay in both semantic dose and
total work, so this result qualifies replay for integration on this bank but
does not isolate pure replay causality. It establishes neither independent-bank
generalization, connected use, native action improvement, parenting, nor a
lifetime effect.

The authoritative result audit is
\path{research_notes/analysis/2026-09-13_event_retention_v2_three_seed_final_fresh_audit.md}
(commit \texttt{d7111125}). The final evidence archive has SHA256
\path{e0492b7ec5c5034490ee19c848a36b40884a780d4675ffd4749cca407b7f7438};
exact reproduction additionally requires the four antecedent archives and
source commit enumerated by that audit.

\subsection{Prospective parenting and lifetime results}

\tbd{} Figure 1: the parented Think--Dream--Sleep child and frozen-parameter
active-memory reference,
with the full four-cell held-out value curves and life-level uncertainty.

\tbd{} Figure 2: entry-adjusted parenting-by-learning interaction plus
proposal quality versus free routing, showing that any gain is neither a
starting-score nor interface/action-volume artifact.

\tbd{} Table 1: final/nAUC/early-slope/forward/backward results and resource
accounting for \texttt{R0}, \texttt{U0}, \texttt{P0}, \texttt{U1}, and
\texttt{P1}.

\tbd{} Table 2: adapter-off, shuffled/wrong-child parenting, and conditional
terminal writer/carrier interventions.

\section{Related work}

Batch experience internalization includes Early Experience
\citep{zhang2025earlyexperience} and LEAFE \citep{ge2026leafe}; the
latter reflects on failures, rolls back, generates improved branches, and
distills those actions into the acting model while removing the explicit
reflection at test time. Active external-memory systems such as MemoPilot
\citep{cai2026memopilot} and Evo-Memory \citep{wei2025evomemory} instead
refine textual state while freezing the actor. Our intended
distinction is repeated, personal, within-life parametric writes and their
online stability, not the first internalization of reflective experience.

\begin{table*}[t]
\centering
\small
\begin{tabular}{lllll}
\toprule
Method family & Learned state & Experience source & Update schedule & Acting policy at deployment \\
\midrule
Early Experience \citep{zhang2025earlyexperience} & full model & alternative actions/future states & batch mid-training & updated once \\
LEAFE \citep{ge2026leafe} & full model & reflection and recovered branches & batch post-training & updated once \\
MemoPilot \citep{cai2026memopilot} & textual memory policy & streaming task trajectories & online text update & frozen \\
Evo-Memory/ReMem \citep{wei2025evomemory} & external textual state & streaming action/thought traces & online text update & frozen \\
Experience Models (target) & per-life low-rank state & own typed actions/outcomes/corrections & periodic within life & repeatedly updated \\
\bottomrule
\end{tabular}
\caption{Mechanism-level positioning. The final row is the setting tested in
this paper, not a claim that its benefit is already established.}
\label{tab:positioning}
\end{table*}

\tbd{} Complete primary-source citations and distinguish the exact paper
implementations from the method-family summary in Table~\ref{tab:positioning}.

\section{Limitations and conclusion}

Per-life adaptation can amplify self-generated errors, erase interfaces,
overfit a narrow environment, and forget prior competence. A compiler gym does
not by itself establish general learning-to-learn, creativity, or physical
world modeling. Parenting is useful only if its effect survives parent removal
and transfers beyond taught content. We therefore separate the narrow
registered claims from the longer-term experience-model architecture.

\section*{Reproducibility statement}
\tbd{} Reference the immutable manifests, source and artifact hashes, exact
commands, seeds, data partitions, acceptance gates, analysis scripts, and
anonymous code supplement.

\section*{AI use statement}
Generative AI systems assisted with implementation, literature search,
experimental-design critique, artifact analysis, result interpretation, and
drafting support. The human author originated and directed the research
program, made the final architecture, protocol, and claim decisions, supplied
extensive iterative scientific feedback, audited the evidence, and wrote or
substantively revised the final manuscript. AI outputs were treated as
proposals rather than evidence and were accepted only after artifact-level or
primary-source verification. We retained hash-bound deliberation and audit
records and did not promote exploratory runs with unresolved provenance or
causal limitations into headline evidence. The human author reviewed the final
code, experiments, numerical claims, citations, and manuscript and takes
responsibility for the submission. \tbd{} Add exact systems/versions and final
human verification assignments from \texttt{AI_USE_LOG.md}; revise singular or
plural authorship after the author list is frozen.

\bibliography{references}
\bibliographystyle{iclr2027_conference}

\appendix
\section{Extended contracts and diagnostics}
\tbd{} Provenance contract, writer fixtures, complete failure table, prompt
and tool schemas, additional curves, and all exploratory-result boundaries.

\end{document}
